\documentclass[11pt]{article}

\usepackage[a4paper,margin=1in]{geometry}
\usepackage[T1]{fontenc}
\usepackage{lmodern}
\usepackage{graphicx}
\usepackage{amsmath,amssymb}
\usepackage{booktabs}
\usepackage{array}
\usepackage{tabularx}
\usepackage{longtable}
\usepackage{float}
\usepackage{subcaption}
\usepackage{hyperref}
\usepackage{enumitem}
\usepackage{xcolor}
\usepackage{microtype}
\usepackage{parskip}

\hypersetup{
    colorlinks=true,
    linkcolor=blue,
    urlcolor=blue,
    citecolor=blue
}

\setlength{\parskip}{6pt}

\title{%
  \textbf{Wafer Map Anomaly Detection on WM-811K}\\[0.4em]
  \large From Reconstruction Baselines to Pretrained Local Feature Modelling
}
\author{Deep Learning Project Report, Group 08\\[0.3em]
  Henry Lee Jun (1004219) \quad Chia Tang (1007200) \quad Genson Low (1005931)}
\date{\today}

\begin{document}

\begin{titlepage}
  \centering
  \vspace*{2cm}
  {\LARGE\bfseries Wafer Map Anomaly Detection on WM-811K\par}
  \vspace{0.5cm}
  {\large From Reconstruction Baselines to Pretrained Local Feature Modelling\par}
  \vspace{1.5cm}
  {\large\bfseries Deep Learning Project Report\par}
  \vspace{0.4cm}
  {\large 50.039 Deep Learning, Y2026\par}
  \vspace{0.4cm}
  {\large Group 08\par}
  \vspace{0.6cm}
  {\normalsize
    Henry Lee Jun (1004219) \quad
    Chia Tang (1007200) \quad
    Genson Low (1005931)\par}
  \vspace{0.8cm}

  {\large \today\par}
  \vfill
  {\small All models implemented in PyTorch. This report documents the experimental design,
  results, and diagnostic analysis for normal-only wafer anomaly detection on the
  WM-811K / LSWMD dataset.\par}
\end{titlepage}

\pagenumbering{roman}

\begin{abstract}
This report presents an anomaly-detection study on the WM-811K / LSWMD wafer-map dataset
under a strict normal-only training regime. The project begins with convolutional autoencoders
to establish a stable baseline pipeline, then moves progressively toward pretrained
ImageNet representations, teacher-student distillation, and PatchCore-style local anomaly
scoring. The main finding is that stronger frozen backbones combined with local feature
modelling improve performance substantially over plain reconstruction or global
embedding-distance baselines. The best completed run on the main benchmark is a
direct-$224\!\times\!224$ PatchCore model built on a pretrained ViT-B/16 backbone
(block 6 features, top-10\% patch reduction), achieving $\mathrm{F1}=0.595$,
$\mathrm{AUROC}=0.956$, and $\mathrm{AUPRC}=0.671$ under a deployment-style
95th-percentile validation threshold.
Score histograms and UMAP diagnostics show that the remaining bottleneck is
score-distribution overlap and naive thresholding rather than feature quality alone.
\end{abstract}

\newpage
\tableofcontents
\newpage

\pagenumbering{arabic}

\section{Introduction and Motivation}
\label{sec:problem}

After fabrication, each die on a semiconductor wafer is electrically tested and
assigned a pass or fail label at its grid location, producing a \emph{wafer bin
map} (WBM). These spatial maps are diagnostically valuable because different
failure patterns indicate different root causes: a ring of failures near the wafer
perimeter suggests a different process error than a central cluster or a diagonal
scratch. Automatically detecting such anomalies enables faster wafer disposition
and improves upstream process control.

Supervised multi-class classification is a natural first approach, and prior work
on WM-811K demonstrates that it can be highly effective in closed benchmarks.
Recent CNN and ensemble classifiers report classification accuracies exceeding
99\% on balanced splits drawn from the same eight labelled defect categories
\cite{khatun2026cbam, taha2025survey}. However, these numbers measure a
different task: \emph{``given a wafer known to contain one of these eight defect
types, which one is it?''} Production anomaly detection asks instead:
\emph{``given a randomly drawn wafer from the line, is it defective at all?''}
These are fundamentally different problems. The high benchmark accuracy arises
because the classifier operates only on confirmed defective wafers drawn from
known categories; in production the vast majority of wafers are normal, and
entirely new failure modes can appear without warning.

Beyond this task mismatch, supervised classification faces three compounding
practical limitations in real fabrication environments:

\begin{itemize}[leftmargin=1.5em]
  \item \textbf{Closed-world constraint.} A classifier with 99\% accuracy on
        eight known defect types has exactly 0\% recall on any new failure mode
        introduced by a process change or new equipment. The defect vocabulary
        in a real fab is neither fixed nor complete.
  \item \textbf{Severe label scarcity.} Only $3.1\%$ of WM-811K wafers carry a
        confirmed defect label and $78.7\%$ are entirely unlabelled. Labelling
        a defect requires expert engineers to correlate wafer maps with equipment
        logs, typically days or weeks after the wafer leaves the line.
  \item \textbf{Scaling cost.} Adding a new defect class requires repeating the
        full label-acquisition cycle from scratch. A normal-only system needs
        only verified-clean wafers, which are abundant and cheap to collect.
\end{itemize}

This project therefore adopts \emph{normal-only anomaly detection}: the model
is trained exclusively on the 147,431 labelled-normal wafers. Any wafer that
deviates from the learned normal representation receives a high anomaly score,
regardless of which defect type caused the deviation. No defect labels are
required at any stage: training, threshold selection, or deployment. The
deployment threshold is the 95th percentile of validation-normal scores,
matching the production constraint that only normal wafers are available for
calibration. Because false negatives (missed defects) are the more critical
error in semiconductor inspection, this report prioritises F1 at the deployed
threshold rather than AUROC alone.

\section{Dataset and Evaluation Protocol}

\subsection{Dataset}

The WM-811K dataset contains 811,457 wafer bin maps (Table~\ref{tab:dataset_distribution}).
Each map is converted to a normalised float array in $[0,1]$ and resized using
nearest-neighbour interpolation. Only explicitly labelled-normal wafers are used
for training; all labelled defect types are treated as anomaly at evaluation time.
The unlabelled $78.7\%$ is excluded entirely to prevent label contamination.

\begin{table}[H]
\centering
\caption{WM-811K dataset composition.}
\label{tab:dataset_distribution}
\begin{tabular}{lcc}
\toprule
Category & Count & Fraction \\
\midrule
No-label & ${\approx}638{,}507$ & $78.7\%$ \\
Labelled: None (Normal) & $147{,}431$ & $18.2\%$ \\
Labelled: Pattern (Defect) & $25{,}519$ & $3.1\%$ \\
\midrule
\textbf{Total} & \textbf{811,457} & \textbf{100\%} \\
\bottomrule
\end{tabular}
\end{table}

Most experiments resize wafer maps to $64\!\times\!64$ pixels. Later
pretrained-backbone experiments additionally evaluate direct $224\!\times\!224$
preprocessing to match each backbone's native resolution.

\subsection{Benchmark Split}

All experiments share a single split drawing 50,000 labelled-normal wafers,
with defects held out exclusively in the test partition (Table~\ref{tab:main_split}).
Normals are divided 80/10/10 across train, validation, and test.
Defects are capped at 5\% of test-normal count to reflect realistic production
prevalence. The validation set serves one purpose only: deriving the deployment
threshold; no defect labels are visible before test evaluation.

\begin{table}[H]
\centering
\caption{Main benchmark split.}
\label{tab:main_split}
\begin{tabular}{lcc}
\toprule
Subset & Normal wafers & Defect wafers \\
\midrule
Train   & 40{,}000 & 0   \\
Val     & 5{,}000  & 0   \\
Test    & 5{,}000  & 250 \\
\bottomrule
\end{tabular}
\end{table}

\subsection{Defect Categories}

The 250 test-set defects span eight pattern types that vary substantially in
spatial scale and coverage (Table~\ref{tab:defect_types}, Figure~\ref{fig:defect_gallery}).

\begin{table}[H]
\centering
\caption{Defect-type breakdown of the main benchmark test set.}
\label{tab:defect_types}
\begin{tabularx}{\linewidth}{lc>{\raggedright\arraybackslash}X}
\toprule
Defect type & Count & Spatial character \\
\midrule
Edge-Ring & 84  & Continuous ring of failures near the wafer perimeter \\
Edge-Loc  & 53  & Localised cluster at or near the wafer edge \\
Center    & 50  & Compact cluster at the wafer centre \\
Loc       & 34  & Off-centre localised patch \\
Scratch   & 15  & Thin elongated streak across the die grid \\
Donut     & 7   & Curved partial ring inside the wafer body \\
Random    & 5   & Scattered failures with no spatial coherence \\
Near-full & 2   & Defective region covering nearly the entire wafer \\
\bottomrule
\end{tabularx}
\end{table}

\begin{figure}[H]
    \centering
    \includegraphics[width=\linewidth]{images/defect_type_gallery.png}
    \caption{Representative examples of all eight defect pattern types in the
    WM-811K evaluation set. Each image is a wafer bin map rendered as a
    binary die-level pass/fail grid. The diversity in spatial scale and
    structure, from full-wafer ring patterns to single-streak scratches,
    is a primary challenge for any global anomaly scoring approach.}
    \label{fig:defect_gallery}
\end{figure}

Broad defects such as Edge-Ring and Center are detectable by any model that
measures overall reconstruction quality. Small, localised defects such as
Scratch and Loc require spatially precise, patch-level scoring. This
distinction drives the experimental progression.

\section{Methods}
\label{sec:methodology}

Table~\ref{tab:method_overview} summarises all model families. The progression
moves from fully in-domain reconstruction models through frozen pretrained
backbones to local patch-level scoring, testing four questions along the way.

\begin{table}[H]
\centering
\caption{Overview of model families evaluated in this study.}
\label{tab:method_overview}
\begin{tabularx}{\linewidth}{>{\raggedright\arraybackslash}p{3.2cm}
                              >{\raggedright\arraybackslash}p{2.2cm}
                              >{\raggedright\arraybackslash}p{2.4cm}
                              >{\raggedright\arraybackslash}X}
\toprule
Method & Training & Scoring & Backbones / notes \\
\midrule
Convolutional Autoencoder
  & From scratch
  & Pixel reconstruction error
  & In-domain encoder, multiple architectural variants \\
Variational Autoencoder (VAE)
  & From scratch
  & Reconstruction error + KL divergence
  & KL weight $\beta$ swept \\
Deep SVDD
  & From scratch
  & Distance to latent hypersphere centre
  & One-class geometric objective, no reconstruction \\
\midrule
Pretrained backbone (center-distance)
  & None (frozen)
  & Global embedding distance to normal centroid
  & ResNet18, WideResNet50-2 \\
\midrule
Teacher-student distillation
  & Student trained on normals
  & Spatial discrepancy map (patch-level)
  & ResNet18, ResNet50, WideResNet50-2 \\
PatchCore
  & None (memory bank only)
  & Patch nearest-neighbour distance
  & ResNet18, ResNet50, WideResNet50-2, EfficientNet-B0, ViT-B/16 \\
FastFlow
  & Normalising flow trained on normals
  & Log-likelihood under learned feature distribution
  & WideResNet50-2 (multilayer) \\
\bottomrule
\end{tabularx}
\end{table}

\subsection{Evaluation Protocol}

Every model uses the same deployment threshold derived from validation-normal
scores alone, with no access to defect labels:
\[
  \tau_{0.95} = \text{95th percentile of validation-normal anomaly scores.}
\]

\begin{table}[H]
\centering
\caption{Metrics reported for each model. Val-threshold F1 is the primary metric.}
\label{tab:metrics}
\begin{tabularx}{\linewidth}{>{\raggedright\arraybackslash}p{3.2cm}
                              >{\raggedright\arraybackslash}X
                              >{\raggedright\arraybackslash}p{2cm}}
\toprule
Metric & Description & Role \\
\midrule
\textbf{Val-threshold F1} & Harmonic mean of precision and recall at $\tau_{0.95}$ & \textbf{Primary} \\
Precision @ $\tau_{0.95}$ & Fraction of flagged wafers that are truly defective & Supporting \\
Recall @ $\tau_{0.95}$ & Fraction of defective wafers successfully detected & Supporting \\
AUROC & Ranking quality across all thresholds & Ranking \\
AUPRC & Ranking quality under class imbalance (more informative than AUROC at 5\% defect rate) & Ranking \\
Best-sweep F1 & Best achievable F1 by sweeping thresholds on test labels & Ceiling only \\
\bottomrule
\end{tabularx}
\end{table}

Best-sweep F1 is never used as a deployment result; it quantifies how much
headroom remains if a more principled threshold rule were available. All models
are implemented in PyTorch (\texttt{torch}, \texttt{torchvision}, \texttt{timm},
\texttt{scikit-learn}, \texttt{umap-learn}; full versions in \texttt{README.md}).

\subsection{Model Families}

The experiments answer four questions in sequence. Full training configurations
and hyperparameter details are listed in Appendix~\ref{app:notebooks}.

\textbf{Q1: Can reconstruction error detect anomalies?} The convolutional
autoencoder (latent dim 128, Adam $10^{-3}$, early stopping) establishes the
baseline pipeline and exposes the key weakness: broad defects such as Edge-Ring
are detectable, but small localised defects (Scratch, Loc) are largely missed
by full-image reconstruction error. A family of variants (BatchNorm, dropout
sweep $\{0.00, 0.05, 0.10, 0.20\}$, residual skip connections, $128\!\times\!128$
resolution, and \emph{$224\!\times\!224$ at pretrained-model native resolution})
tests whether this ceiling can be raised within the reconstruction paradigm.
A VAE ($\beta$ swept) and Deep SVDD (one-class hypersphere) are also included
as alternative one-class baselines. A post-training score ablation across seven
reduction rules (MSE, MAE, max pixel, top-$k$ mean, foreground variants) reveals
that scoring choice is as impactful as architecture.

\textbf{Q2: Do pretrained features help without local scoring?} Frozen
ImageNet-pretrained ResNet18 and WideResNet50-2 are evaluated using global
center-distance scoring (L2 from train-normal feature centroid). This
deliberately stress-tests whether rich backbone features alone are sufficient
before any spatial scoring is introduced.

\textbf{Q3: Does local spatial scoring matter?} Three methods test patch-level
anomaly signals. \emph{Teacher-student distillation} trains a lightweight student
CNN to reproduce frozen teacher spatial feature maps on normal wafers; the
per-location mismatch map is reduced to a wafer score by averaging the most
deviant patch locations (top-k ratio swept; branch weights swept post-training).
Backbones: ResNet18, ResNet50, WideResNet50-2 (single-layer and multilayer).
\emph{PatchCore} stores a coreset memory bank of patch embeddings from one
forward pass, then scores test wafers by nearest-neighbour patch distance,
training-free. Backbones: ResNet18, ResNet50, WideResNet50-2, EfficientNet-B0.
\emph{FastFlow} fits a normalising flow on frozen WideResNet50-2 feature maps,
scoring by log-likelihood under the learned feature distribution.

\textbf{Q4: Does input resolution matter?} Wafer maps are typically
$32$--$64$ pixels; pretrained backbones expect $224\!\times\!224$. In the
$64\!\times\!64$ experiments maps are upsampled inside the model, potentially
discarding spatial detail. This question is evaluated by caching images directly
at the backbone's native resolution (EfficientNet-B0: $224\!\times\!224$;
EfficientNet-B1: $240\!\times\!240$; WideResNet50-2: $224\!\times\!224$;
ViT-B/16: $224\!\times\!224$), each compared against its $64\!\times\!64$
counterpart with all other settings fixed. To ensure a fair reconstruction-based
baseline, a plain convolutional autoencoder is also trained at the native
$224\!\times\!224$ resolution for direct comparison against pretrained methods.
ViT-B/16 is also qualitatively distinct: it processes the image as a sequence of
non-overlapping $16\!\times\!16$ patches with global self-attention, producing
context-aware patch features from the first layer onwards.

\section{Main Results}

\subsection{Overview}

Table~\ref{tab:main_leaderboard} summarises ten representative models
that anchor the main narrative. Full sweep results across all variants are given in
Appendix~\ref{app:notebooks}.

\begin{table}[H]
\centering
\caption{Representative results on the main benchmark (ranked by val-threshold F1).
Models are selected to illustrate each stage of the experimental progression
rather than to list every variant tested.}
\label{tab:main_leaderboard}
\begin{tabularx}{\linewidth}{>{\raggedright\arraybackslash}Xcccc}
\toprule
Model & F1 & AUROC & AUPRC & Sweep F1 \\
\midrule
PatchCore + ViT-B/16, direct $224\!\times\!224$              & \textbf{0.595} & \textbf{0.956} & \textbf{0.671} & \textbf{0.651} \\
PatchCore + EfficientNet-B1, direct $240\!\times\!240$       & 0.591 & 0.935 & 0.609 & 0.651 \\
PatchCore + WideResNet50-2, direct $224\!\times\!224$        & 0.549 & 0.931 & 0.659 & 0.634 \\
PatchCore + WideResNet50-2, $64\!\times\!64$                 & 0.532 & 0.917 & 0.562 & 0.586 \\
Teacher-Student + ResNet50                                   & 0.525 & 0.909 & 0.599 & 0.606 \\
Teacher-Student + WideResNet50-2 (layer2+layer3)             & 0.524 & 0.923 & 0.546 & 0.561 \\
Autoencoder (plain, $224\!\times\!224$, native resolution)   & 0.510 & 0.901 & 0.596 & 0.587 \\
Autoencoder + BatchNorm (best score)                         & 0.502 & 0.834 & 0.568 & 0.630 \\
FastFlow + WideResNet50-2                                    & 0.482 & 0.871 & 0.489 & 0.482 \\
WideResNet50-2 center-distance (frozen)                      & 0.243 & 0.677 & 0.142 & 0.270 \\
ResNet18 center-distance (frozen)                            & 0.236 & 0.685 & 0.195 & 0.260 \\
\bottomrule
\end{tabularx}
\end{table}

\begin{figure}[H]
    \centering
    \includegraphics[width=\linewidth]{images/overall_experiment_comparison.png}
    \caption{Overall comparison across all completed experiments. Left: top runs
    ranked by deployed F1. Right: all runs in AUPRC-vs-F1 space, point size
    scaled by AUROC. The gap between the reconstruction family (lower left) and
    the local pretrained-feature methods (upper right) is the central finding
    of the project.}
    \label{fig:overall_comparison}
\end{figure}

Three cross-cutting findings run through every model family and explain the
shape of the leaderboard.

\textbf{1. Score design matters almost as much as model design.}
The BatchNorm autoencoder with its default scoring rule ranked \emph{below}
the plain baseline ($\mathrm{F1}=0.431$ vs $0.468$). Post-training rescoring
with maximum absolute pixel error lifted it to $\mathrm{F1}=0.502$ without
any retraining. The same interaction appeared in the teacher-student family:
the ResNet50 teacher at its default score reached $\mathrm{F1}=0.488$. A
post-training score sweep lifted it to $\mathrm{F1}=0.525$. The scoring rule
is a design decision co-equal with the model architecture.

\textbf{2. Pretrained backbones require local scoring to be useful.}
Both center-distance baselines ($\mathrm{F1}<0.25$, $\mathrm{AUROC}<0.70$)
are weaker than every reconstruction model despite using far richer feature
representations. The bottleneck is not feature quality but global aggregation:
averaging all spatial information into one embedding vector discards exactly
the local patch evidence that defects produce.

\textbf{3. Native high-resolution preprocessing provides a clear and consistent lift.}
EfficientNet-B0 PatchCore rose from $\mathrm{F1}=0.467$ to $0.544$ when
wafer images were cached at full resolution rather than upsampled from
$64\!\times\!64$ inside the model. Scaling backbone capacity from B0 to B1 at
its native $240\!\times\!240$ input resolution pushes performance further to
$\mathrm{F1}=0.591$, making it the strongest CNN PatchCore result in the project.
WideResNet50-2 PatchCore improved similarly at $224\!\times\!224$, and
the ViT-B/16 branch achieves the overall best result ($\mathrm{F1}=0.595$)
at that same resolution.

\subsection{Reconstruction Family}

\begin{figure}[H]
\centering
\includegraphics[width=0.95\linewidth]{images/autoencoder/ae64_training_and_evaluation.png}
\caption{Baseline AE: training curves and score distribution.}
\label{fig:ae64_training}
\end{figure}

\begin{figure}[H]
\centering
\includegraphics[width=0.95\linewidth]{images/autoencoder/ae64_score_ablation.png}
\caption{Score ablation on the baseline AE checkpoint. The same checkpoint spans a wide range of F1 values depending solely on which score is applied.}
\label{fig:ae64_score_ablation}
\end{figure}

The plain autoencoder, scored with a top-$k$ percentile mean, achieves
$\mathrm{F1}=0.468$, $\mathrm{AUROC}=0.839$, $\mathrm{AUPRC}=0.522$.
The per-class pattern is informative: broad, globally visible defects such
as Edge-Ring (recall 0.810) and Center (recall 0.720) are well-detected,
while small, localised defects such as Loc (recall 0.206) and Scratch
(recall 0.200) are largely missed. This failure pattern motivates every
subsequent modelling direction.


\begin{figure}[H]
\centering
\begin{subfigure}[t]{\linewidth}
  \includegraphics[width=0.85\linewidth]{images/autoencoder/reconstruction_examples.png}
  \caption{Reconstruction mechanism: inputs (top) and reconstructions (bottom) for 3 normal and 3 anomalous wafers.}
\end{subfigure}
\vspace{0.3cm}
\begin{subfigure}[t]{\linewidth}
  \includegraphics[width=0.5\linewidth]{images/autoencoder/failure_examples_fn.png}
  \caption{False negatives: small defects (Loc, Scratch) produce high reconstruction error (visible in right column) yet remain undetected, exposing the scoring rule's weakness on sparse patterns.}
\end{subfigure}
\caption{Baseline autoencoder failure analysis: the reconstruction mechanism (top) successfully identifies defects through elevated error, but the scoring rule (top-$k$ mean) underweights small, sparse defects, causing missed detections (bottom) despite strong error signals.}
\label{fig:ae_reconstruction_failure}
\end{figure}

The root cause is not the reconstruction mechanism but the \emph{scoring rule}. Adding batch normalisation with the same top-$k$ mean scoring rule counterintuitively \emph{reduces} F1 to 0.431. However, score ablation reveals the opportunity: batch normalisation changes the error distribution so that the single worst pixel becomes a more reliable anomaly signal than averaging. Rescoring the BatchNorm checkpoint with maximum absolute pixel error lifts F1 to 0.502, recall to 0.668, and AUPRC to 0.568, the best reconstruction-family result in the project. This is achieved without retraining.

\textbf{Resolution Impact:} A critical question is whether the $64\!\times\!64$ resolution limits reconstruction-based detection. A plain autoencoder trained at the native $224\!\times\!224$ resolution (the intended input size for pretrained ImageNet backbones) achieves $\mathrm{F1}=0.510$, $\mathrm{AUROC}=0.901$, and $\mathrm{AUPRC}=0.596$. The ranking quality improvement is dramatic: AUROC jumped $+0.067$ percentage points without any architectural changes. This establishes a fair baseline for comparing reconstruction-based and pretrained-feature methods at their native resolutions.

\begin{figure}[H]
\centering
\includegraphics[width=0.95\linewidth]{experiments/anomaly_detection/autoencoder/x224/main/artifacts/autoencoder_x224/plots/score_ablation_summary.png}
\caption{Score ablation on the x224 autoencoder checkpoint. Seven scoring rules (MSE, MAE, max pixel, top-$k$ mean, foreground variants, pooled MAE) are evaluated on the same checkpoint, revealing that top-$k$ mean continues to be the dominant scoring strategy at higher resolution.}
\label{fig:ae224_score_ablation}
\end{figure}

\begin{figure}[H]
\centering
\includegraphics[width=\linewidth]{images/autoencoder/autoencoder_family_comparison.png}
\caption{Autoencoder family: all AE-family variants ranked by deployed F1 (left)
and precision-recall operating points (right). The $224\!\times\!224$ variant (top-$k$ mean)
now ranks first in the family, demonstrating that resolution improvement dominates
architectural changes like BatchNorm and residual skip connections.}
\label{fig:ae_family}
\end{figure}

Two alternative one-class approaches were also evaluated as comparisons.
The Variational Autoencoder, with KL weight swept over four values, peaked
at $\mathrm{F1}=0.340$ ($\beta=0.005$), remaining clearly below the plain
autoencoder despite the more principled probabilistic latent space. Deep
SVDD, which maps normal wafers into a compact hypersphere rather than
performing reconstruction, reached $\mathrm{F1}=0.360$ and
$\mathrm{AUROC}=0.788$, but with very weak $\mathrm{AUPRC}=0.213$,
indicating poor ranking quality under class imbalance. Neither approach
is competitive with the tuned autoencoder. Additional reconstruction
variants including dropout sweeps, a residual encoder-decoder, and a
$128\!\times\!128$ version are documented in the appendix. None improved
on the BatchNorm result.

\subsection{Pretrained Backbone Baselines}

\begin{figure}[H]
\centering
\includegraphics[width=\linewidth]{images/compact_baseline_comparison.png}
\caption{Non-leading baselines: VAE, Deep SVDD, and plain embedding-distance
backbones. All fall well below the reconstruction-family ceiling, establishing
that both probabilistic one-class methods and raw pretrained features are
insufficient without local scoring.}
\label{fig:compact_baselines}
\end{figure}

Frozen ResNet18 ($\mathrm{F1}=0.236$, $\mathrm{AUROC}=0.685$) and
WideResNet50-2 ($\mathrm{F1}=0.243$, $\mathrm{AUROC}=0.677$) with
center-distance scoring both fall \emph{below} the plain autoencoder
despite using much richer ImageNet features. The failure is not the
backbone but the aggregation: collapsing a full spatial feature map into
a single global embedding vector loses exactly the local deviation signals
that defects produce. Larger backbones do not fix this. Only local scoring
methods do. These results are therefore deliberate negative controls that
motivate the local scoring methods in the next two subsections.

\subsection{Teacher-Student Distillation}

\begin{figure}[H]
\centering
\includegraphics[width=\linewidth]{images/teacherstudent/ts_family_comparison_full.png}
\caption{Teacher-student family: ResNet18 vs ResNet50 teachers. Left: deployed
F1, AUPRC, and AUROC. Right: precision-recall operating points. The ResNet50
teacher clearly separates from the ResNet18 teacher after score selection.}
\label{fig:ts_family}
\end{figure}

Teacher-student distillation addresses the local scoring problem by training
a lightweight student network to reproduce the teacher's spatial feature maps
on normal wafers only. At inference, the per-location discrepancy between
teacher and student forms a spatial anomaly map, which is then reduced to a
wafer score by averaging over the most deviant patch locations.

The ResNet18 teacher with the default combined score is weak, but rescoring
with student-only top-$k$ mean at ratio 0.20 lifts F1 to 0.495. The pattern
repeats with the ResNet50 teacher: the default score gives $\mathrm{F1}=0.488$,
while a post-training branch-weight sweep yields the best teacher-student result
at $\mathrm{F1}=0.525$, $\mathrm{AUROC}=0.909$, $\mathrm{AUPRC}=0.599$.
This is the strongest result in the project at this stage and confirms
that local spatial scoring on pretrained features decisively outperforms both
reconstruction models and global embedding distance.

Scaling the teacher backbone further to WideResNet50-2 was the natural next
step. Using a single teacher feature layer (layer2) gives $\mathrm{F1}=0.508$,
$\mathrm{AUROC}=0.920$. Combining two spatial scales (layer2 and layer3)
gives $\mathrm{F1}=0.524$, $\mathrm{AUROC}=0.923$, $\mathrm{AUPRC}=0.546$,
comparable to the ResNet50 teacher on deployed F1 but with clearly stronger
anomaly ranking quality. These experiments validate WideResNet50-2 as a strong
feature source for local anomaly detection and directly motivate its subsequent
use as a PatchCore backbone. Full sweep details are in the appendix.

\begin{figure}[H]
\centering
\includegraphics[width=\linewidth]{images/teacherstudent/ts_family_comparison_full.png}
\caption{Teacher-student distillation: all four backbone variants compared across
F1, AUROC, and AUPRC. WideResNet50-2 multilayer achieves the strongest AUROC
(0.923) among the teacher-student variants, confirming it as the best teacher
backbone. F1 is near-identical between ResNet50 and WRN50-2 multilayer (0.525
vs 0.524), but WRN50-2 provides clearly stronger anomaly ranking quality.
Generated by \texttt{experiments/anomaly\_detection/report\_figures/notebook.ipynb}.}
\label{fig:ts_family_full}
\end{figure}

Teacher-student distillation has one persistent drawback: it requires training
a student network for 25--30 epochs and carefully tuning the branch weighting
between the student and auxiliary autoencoder components after training. The
next natural question is whether the spatial discrepancy signal can be recovered
more directly and without any training at all.

\subsection{FastFlow: Flow-Based Feature Density}

FastFlow tests an alternative local density approach: rather than training a
student to mimic teacher features, it fits a normalising flow directly on top
of frozen WideResNet50-2 spatial feature maps, assigning anomaly scores through
log-likelihood under the learned normal feature distribution. This requires
training the flow head but not the backbone, and unlike PatchCore it does not
need a stored memory bank at inference.

The selected multilayer configuration (layer2 + layer3, 4 flow steps, plain
mean reduction) achieves $\mathrm{F1}=0.482$, $\mathrm{AUROC}=0.871$,
$\mathrm{AUPRC}=0.489$. This is competitive with the reconstruction-family
ceiling ($\mathrm{F1}=0.502$) and demonstrates that flow-based density
modelling on pretrained features is a viable approach. However, it falls
clearly below the best teacher-student result ($\mathrm{F1}=0.525$) and
even further below the best WideResNet50-2 PatchCore result ($\mathrm{F1}=0.532$).
Per-defect recall shows a distinctive pattern: FastFlow is relatively strong
on medium-scale localised defects (Loc recall 0.588, Edge-Loc 0.528) but
especially weak on the thin Scratch class (recall 0.133), suggesting the
learned density boundary does not capture the fine spatial geometry of narrow
streak defects. The result positions FastFlow as a meaningful reference point
but confirms that nearest-neighbour memory-bank scoring (PatchCore) is a
stronger approach on this dataset when the same backbone is used.

\begin{figure}[H]
\centering
\begin{subfigure}[t]{0.48\linewidth}
  \includegraphics[width=\linewidth]{images/fastflow/variant_comparison_metrics.png}
  \caption{Variant comparison: F1, AUROC, AUPRC across three configurations.}
\end{subfigure}
\hfill
\begin{subfigure}[t]{0.48\linewidth}
  \includegraphics[width=\linewidth]{images/fastflow/best_variant_defect_breakdown.png}
  \caption{Per-defect recall for the selected variant (\texttt{wrn50\_l23\_s4}, mean reduction).}
\end{subfigure}
\caption{FastFlow results. Left: the multilayer $4$-step variant (layer2+layer3)
outperforms both the 6-step multilayer and the single-layer configurations,
showing that feature breadth matters more than flow depth. Right: per-defect
recall pattern. Loc and Edge-Loc are the relative strengths, while Scratch (0.133)
is the clear weak point, consistent with the broader project finding that
thin linear defects are the hardest for all feature-density methods.}
\label{fig:fastflow_results}
\end{figure}

\subsection{PatchCore: Local Nearest-Neighbour Scoring}

\begin{figure}[H]
\centering
\includegraphics[width=\linewidth]{images/patchcore/patchcore_family_comparison.png}
\caption{PatchCore family: best result per backbone (left) and all PatchCore
variants coloured by backbone (right). The progression from in-domain AE
encoder to pretrained ResNet to WideResNet50-2 to ViT-B/16 follows a clear
upward trajectory, with native $224\!\times\!224$ preprocessing providing the
final decisive lift.}
\label{fig:patchcore_family}
\end{figure}

PatchCore offers a training-free alternative to teacher-student distillation:
rather than training a student network, it stores a memory bank of patch-level
feature embeddings extracted from normal wafers in a single forward pass, then
scores test wafers by the nearest-neighbour distance of their patch embeddings
to the bank. Because each patch is scored independently, the method is
inherently sensitive to localised defects, and it requires no additional
optimisation loop once the backbone is chosen.

The backbone progression tells a clear story of escalating feature quality.
As an initial sanity check, the frozen BatchNorm AE encoder (the same checkpoint
used in the reconstruction family) was used as the PatchCore backbone. This
produced $\mathrm{F1}=0.336$, immediately confirming that local patch scoring
over in-domain features is useful even though the features themselves are weak.
Switching to ImageNet-pretrained ResNet18 jumps to $\mathrm{F1}=0.401$,
validating that pretrained features are the right direction. ResNet50 improves
further to $\mathrm{F1}=0.420$. Scaling to the multilayer WideResNet50-2
backbone at $64\!\times\!64$ reaches $\mathrm{F1}=0.532$, which for the first
time exceeds both the best autoencoder and the best teacher-student result,
albeit by a modest margin. The decisive separation comes from moving to
native high-resolution preprocessing.
\begin{figure}[H]
\centering
\includegraphics[width=\linewidth]{experiments/anomaly_detection/report_figures/artifacts/plots/patchcore_backbone_progression.png}
\caption{PatchCore backbone progression ranked by val-threshold F1. The dashed
line marks the best reconstruction-family result (AE-BN, $\mathrm{F1}=0.502$).
The large jump at WideResNet50-2 x64 ($+0.112$) reflects the gain from
multilayer pretrained features over single-backbone CNN models. The subsequent
gains from native-resolution preprocessing and backbone scaling are also visible.
Generated by \texttt{experiments/anomaly\_detection/report\_figures/notebook.ipynb}.}
\label{fig:patchcore_progression}
\end{figure}

The EfficientNet-B0 branch provides the clearest controlled test of the
resolution hypothesis: the same backbone and method at $64\!\times\!64$
preprocessing gives $\mathrm{F1}=0.467$, while switching to native
$224\!\times\!224$ preprocessing raises it to $\mathrm{F1}=0.544$. The spatial
information that was being discarded by upsampling matters substantially.
WideResNet50-2 PatchCore at native $224\!\times\!224$ reaches $\mathrm{F1}=0.549$,
confirming the pattern.

\begin{figure}[H]
\centering
\includegraphics[width=0.75\linewidth]{images/resolution_effect_comparison.png}
\caption{Controlled resolution comparison for EfficientNet-B0 and WideResNet50-2
PatchCore. All other settings are held fixed; only the wafer preprocessing
resolution changes. The gain is larger for EfficientNet-B0 ($+0.077$ F1) than
for WideResNet50-2 ($+0.017$ F1), suggesting EfficientNet-B0 is more sensitive
to the spatial detail captured at higher resolution.
Generated by \texttt{experiments/anomaly\_detection/report\_figures/notebook.ipynb}.}
\label{fig:resolution_comparison}
\end{figure}

\begin{figure}[H]
\centering
\includegraphics[width=\linewidth]{images/patchcore/umap_resolution_comparison_effb0.png}
\caption{UMAP embedding geometry for the \emph{same} EfficientNet-B0 backbone and PatchCore scoring at two input resolutions: 64$\times$64 upsampled (top 2 panels, F1$=0.467$) vs.\ native 224$\times$224 (bottom 2 panels, F1$=0.544$). Each resolution shows defect-type coloring (left) and anomaly score heatmap (right). Backbone weights, memory bank size, and scoring rule are identical --- only the cached image resolution differs. At 64$\times$64, defect clusters scatter throughout the normal cloud with no clear boundary. At native 224$\times$224, defect types concentrate in distinct peripheral regions, particularly Edge-Ring and Center. The geometry improvement directly accounts for the $+0.077$ F1 gain without any architecture change. Embeddings extracted from the saved checkpoint and UMAP run in \texttt{experiments/anomaly\_detection/report\_figures/notebook.ipynb}.}
\label{fig:umap_resolution_effb0}
\end{figure}


\subsubsection{EfficientNet-B1 PatchCore: Strongest CNN Result}

Scaling from EfficientNet-B0 to EfficientNet-B1 at its native
$240\!\times\!240$ input resolution, with one-layer feature extraction from
block 3 and a coreset memory bank of 240k patches, pushes the CNN PatchCore
branch to its best result: $\mathrm{F1}=0.591$, $\mathrm{AUROC}=0.935$,
$\mathrm{AUPRC}=0.609$, Sweep F1 $=0.651$. This result is only 0.004 F1
below the ViT-B/16 leader and represents the strongest CNN-based
operating point in the project.

\begin{figure}[H]
\centering
\begin{subfigure}[t]{\linewidth}
  \includegraphics[width=\linewidth]{images/patchcore/effb1/benchmark_distribution_sweep.png}
  \caption{Score distribution and threshold sweep, EfficientNet-B1 $240\!\times\!240$.}
\end{subfigure}
\vspace{0.4cm}
\begin{subfigure}[t]{\linewidth}
  \includegraphics[width=\linewidth]{images/patchcore/effb1/benchmark_defect_breakdown.png}
  \caption{Per-defect recall, EfficientNet-B1 $240\!\times\!240$.}
\end{subfigure}
\caption{EfficientNet-B1 PatchCore at $240\!\times\!240$: score distribution with
deployed threshold and per-defect recall. Edge-Ring, Donut, Random, and Near-full
all reach recall $\geq 0.93$, while Scratch (0.40) and Loc (0.59) remain the
hardest categories, consistent with the broader project pattern.}
\label{fig:effnetb1_x240}
\end{figure}

Per-defect recall for EfficientNet-B1 at $240\!\times\!240$: Scratch 0.400,
Loc 0.588, Center 0.700, Edge-Loc 0.792, Edge-Ring 0.929, Donut 1.000,
Random 1.000, Near-full 1.000. This is competitive with the ViT result on
all categories except Center and Edge-Loc, where the ViT's global context-aware
patch features carry a measurable advantage.

\subsubsection{ViT-B/16 PatchCore: Best Overall Result}

Having exhausted the incremental gains from CNN backbone scaling, the final
upgrade is qualitative: the Vision Transformer (ViT-B/16), which does not use
convolutions at all. Instead of building a spatial hierarchy through strided
layers, ViT divides the input into a fixed grid of $16\!\times\!16$ pixel
patches and processes them as a sequence with multi-head self-attention.
Every patch attends to every other patch at every layer, producing features
that are globally context-aware from the outset, architecturally distinct
from EfficientNet-B1 where receptive fields grow only gradually across layers.
Applied with PatchCore at native $224\!\times\!224$, it achieves the best
result in the project: $\mathrm{F1}=0.595$, $\mathrm{AUROC}=0.956$,
$\mathrm{AUPRC}=0.671$.

\begin{figure}[H]
\centering
\begin{subfigure}[t]{0.55\linewidth}
  \includegraphics[width=\linewidth]{experiments/anomaly_detection/patchcore/wideresnet50/x224/multilayer_umap_followup/artifacts/patchcore-wideresnet50-multilayer-umap/topk_mb50k_r005_x224/plots/score_histogram.png}
  \caption{Score histogram: WideResNet50-2 PatchCore $224\!\times\!224$.}
\end{subfigure}
\hfill
\begin{subfigure}[t]{0.43\linewidth}
  \includegraphics[width=\linewidth]{experiments/anomaly_detection/patchcore/wideresnet50/x224/multilayer_umap_followup/artifacts/patchcore-wideresnet50-multilayer-umap/topk_mb50k_r005_x224/plots/threshold_sweep_metrics.png}
  \caption{Threshold sweep: precision, recall, and F1.}
\end{subfigure}
\caption{Score distribution and threshold sweep for the WideResNet50-2 PatchCore
$224\!\times\!224$ variant. The 95th-percentile validation threshold ($\tau=0.537$)
and the best-sweep threshold ($0.559$) are both marked. The gap between them
represents the headroom that a more principled threshold rule could recover.}
\label{fig:wrn_x224_diag}
\end{figure}

\begin{figure}[H]
\centering
\begin{subfigure}[t]{0.48\linewidth}
  \includegraphics[width=\linewidth]{experiments/anomaly_detection/patchcore/vit_b16/x224/main/artifacts/patchcore_vit_b16_5pct/main_5pct/plots/test_evaluation.png}
  \caption{Score distribution, threshold sweep, and confusion matrix.}
\end{subfigure}
\hfill
\begin{subfigure}[t]{0.48\linewidth}
  \includegraphics[width=\linewidth]{experiments/anomaly_detection/patchcore/vit_b16/x224/main/artifacts/patchcore_vit_b16_5pct/main_5pct/plots/defect_breakdown.png}
  \caption{Per-defect recall.}
\end{subfigure}
\caption{ViT-B/16 PatchCore at direct $224\!\times\!224$, the best result in the
project ($\mathrm{F1}=0.595$, $\mathrm{AUROC}=0.956$, $\mathrm{AUPRC}=0.671$).
Left: score distribution at the 95th-percentile validation threshold and
threshold sweep showing the headroom between the deployed and best-sweep
operating points. Right: per-defect recall showing broad coverage across
all eight defect types, with the hardest categories (Scratch, Loc) now
clearly above 0.7.}
\label{fig:vit_main_result}
\end{figure}

The key per-class gains over EfficientNet-B1 are concentrated in the hardest
localised defects: Scratch improves from 0.400 to 0.727, Loc from 0.588 to
0.829, and Edge-Loc from 0.792 to 0.705 (see Table~\ref{tab:per_defect}).
The ViT's global self-attention mechanism appears to capture the long-range
spatial context of thin streak and cluster patterns that convolutional features
miss. The one category where ViT trails is Center (0.618 vs 0.780 for AE-BN),
a broad concentric defect whose signal is already well captured by global
reconstruction error. This suggests that ViT patch features are calibrated
toward local irregularities rather than global wafer structure.

As a final check, a post-hoc score ensemble combined the WRN50_2 PatchCore
and TS-ResNet50 score streams using normalised-score maximum fusion. This
improved AUPRC from 0.562 to 0.611, reflecting genuine complementarity between
the two scoring methods, but deployed F1 ($0.529$) still fell below the
standalone WRN PatchCore ($0.532$). The ViT result supersedes both by a clear
margin, confirming that backbone quality rather than score fusion is the
higher-leverage lever.

\subsection{Cross-Model Defect Analysis}

Table~\ref{tab:per_defect} and Figure~\ref{fig:per_defect_chart} compare
per-defect recall across six representative models spanning the full
progression, from the reconstruction baseline to the best local feature result.

\begin{table}[H]
\centering
\caption{Per-defect recall for six representative models tracing the experimental
progression. Bold indicates the best recall for that defect type.}
\label{tab:per_defect}
\begin{tabular}{lcccccc}
\toprule
Defect type
  & Baseline AE
  & AE-BN
  & WRN center
  & TS-Res50
  & WRN PC $224$px
  & ViT-B/16 PC \\
\midrule
Scratch   & 0.200 & 0.467 & 0.200 & 0.333 & 0.667 & \textbf{0.727} \\
Loc       & 0.206 & 0.294 & 0.176 & 0.559 & 0.735 & \textbf{0.829} \\
Edge-Loc  & 0.434 & 0.472 & 0.132 & 0.491 & 0.623 & \textbf{0.705} \\
Center    & 0.720 & \textbf{0.800} & 0.140 & 0.720 & 0.780 & 0.618 \\
Donut     & 0.571 & 0.571 & 0.286 & 0.714 & \textbf{1.000} & \textbf{1.000} \\
Edge-Ring & 0.810 & 0.893 & 0.464 & \textbf{0.929} & 0.798 & 0.941 \\
Random    & 0.600 & 0.800 & 0.600 & \textbf{1.000} & \textbf{1.000} & \textbf{1.000} \\
Near-full & 0.500 & \textbf{1.000} & 0.000 & \textbf{1.000} & \textbf{1.000} & \textbf{1.000} \\
\bottomrule
\end{tabular}
\end{table}

Several patterns emerge clearly. The center-distance baseline struggles
uniformly across all defect types, confirming that global feature aggregation
is the bottleneck. The reconstruction family is strongest on Center, a globally
visible defect where full-image reconstruction error carries a clear signal.
Local scoring methods lead on every small, spatially confined defect:
WRN PatchCore at $224\!\times\!224$ and ViT-B/16 PatchCore both achieve
strong results on Scratch, Loc, and Edge-Loc, with ViT-B/16 pulling ahead on
those three categories ($+0.06$ Scratch, $+0.09$ Loc, $+0.08$ Edge-Loc) at
the cost of slightly weaker Center recall (0.618 vs 0.800 for AE-BN).
Teacher-student distillation remains the leader on Edge-Ring (0.929).
Scratch is the hardest category across all models,
reflecting its thin, low-coverage spatial signature.

\begin{figure}[H]
\centering
\includegraphics[width=\linewidth]{experiments/anomaly_detection/report_figures/artifacts/plots/per_defect_model_comparison.png}
\caption{Per-defect recall across the six key models, shown as a grouped bar
chart. The shaded background highlights the three hardest defect classes
(Scratch, Loc, Edge-Loc). The progression from reconstruction-based methods
to local pretrained-feature methods is most visible in those three categories:
Scratch recall rises from 0.20 (Baseline AE) to 0.727 (ViT-B/16 PatchCore),
and Loc recall rises from 0.21 to 0.829. The WRN center-distance baseline
(grey) confirms that pretrained features alone without local scoring perform
worse than any trained model.
Generated by \texttt{experiments/anomaly\_detection/report\_figures/notebook.ipynb}.}
\label{fig:per_defect_chart}
\end{figure}

\begin{figure}[H]
\centering
\includegraphics[width=\linewidth]{experiments/anomaly_detection/report_figures/artifacts/plots/score_distribution_comparison.png}
\caption{Score distributions for normal (blue) and anomaly (red) wafers across
four key models. The 95th-percentile validation threshold is marked with a
dashed line on each panel. As model quality improves, anomaly scores shift
rightward relative to normal scores, but the distributions continue to overlap
near the threshold, explaining both the moderate F1 values and the consistent
gap between deployed and best-sweep F1 across all methods.
Generated by \texttt{experiments/anomaly\_detection/report\_figures/notebook.ipynb}.}
\label{fig:score_distributions}
\end{figure}

\subsection{Comparison with Published Work}

Two published works provide the closest available reference points for
situating these results within the literature on normal-only wafer anomaly
detection. Direct numerical comparison is qualified carefully below, as
evaluation protocols differ across papers.

\paragraph{Jo and Lee (2021): one-class classification on WM-811K.}
The most directly comparable published study is \cite{jo2021oneclass},
which applies one-class classification to the identical WM-811K dataset
under a normal-only training regime. The paper trains and evaluates three
models: Deep SVDD \cite{ruff2018deepsvdd}, a standalone Adversarial
Autoencoder (AAE), and a hybrid AAE with DSVDD prior (AAE-DSVDD). Their
experimental design mirrors the present project in the key respects that
matter for comparison: training uses only labelled-normal wafers, and
evaluation is image-level binary detection across the same eight defect
pattern classes.

Their reported results show that Deep SVDD achieves precision and accuracy
below 0.5 across most classes, with AAE providing moderate improvement and
AAE-DSVDD achieving the best F1 among the three. This trajectory is
consistent with the present work: our Deep SVDD result (F1 = 0.360)
confirms the weakness of the pure hypersphere objective, and our
reconstruction family (best F1 = 0.502) already surpasses the AAE-class
performance reported in that work. The best result in the present project,
ViT-B/16 PatchCore at direct $224\!\times\!224$ (F1 = 0.595,
AUROC = 0.956), substantially exceeds all models reported in
\cite{jo2021oneclass} on image-level ranking quality, extending their
baseline study with the local pretrained-feature methods they do not
consider.

\paragraph{Roth et al.\ (2022): PatchCore on MVTec AD.}
The PatchCore method used as the primary local scoring approach in this
project is introduced by \cite{roth2022patchcore}. In that paper, PatchCore
is evaluated on the MVTec AD benchmark, achieving image-level
AUROC near 99\% across industrial texture and object categories under a
training-free, normal-only regime. While WM-811K is not evaluated in the
original paper, the method transfers directly to this domain: patch-level
nearest-neighbour scoring against a coreset-reduced memory bank of
normal-image features is architecture-agnostic. The present project can
therefore be understood as the first systematic application of PatchCore to
wafer-map anomaly detection under a strict deployment-style threshold
constraint, evaluating five backbone variants across two input resolutions.

The gap between MVTec AD performance (AUROC $\approx$99\%) and the present
WM-811K results (best AUROC = 0.956) is expected and informative. Wafer
maps present a structurally harder problem than most MVTec categories:
defects are sparse, occupying only a small fraction of die area. The normal
wafer texture is nearly uniform, reducing the contrast between normal and
defective patches, and the strict 95th-percentile threshold calibration on
validation normals alone removes the oracle-threshold headroom that
boosts reported AUROC-only benchmarks. The remaining 0.044 AUROC gap and
the deployed F1 of 0.595 are therefore attributable to the domain
difficulty and the deployment threshold constraint rather than to
implementation quality.
\section{Expanded Holdout Validation}

To test whether the strongest checkpointed models remain stable on a much larger
evaluation pool, a disjoint secondary holdout was constructed:

\begin{table}[H]
\centering
\caption{Main benchmark vs.\ expanded holdout dataset sizes.}
\label{tab:holdout_sizes}
\begin{tabular}{lcc}
\toprule
Subset & Main benchmark & Expanded holdout \\
\midrule
Train normals & 40{,}000 & 40{,}000 \\
Val normals   &  5{,}000 &  5{,}000 \\
Test normals  &  5{,}000 & 70{,}000 \\
Test defects  &    250   &  3{,}500 \\
\bottomrule
\end{tabular}
\end{table}

The defect pool expands by roughly 14$\times$, making rare-class recall far more
meaningful than in the main benchmark. The same train and validation normals are
kept fixed, so each reevaluation reuses the previously selected checkpoint and
its validation-derived threshold policy.

\begin{table}[H]
\centering
\caption{Defect-type scaling from benchmark to expanded holdout.}
\label{tab:holdout_defects}
\begin{tabular}{lcc}
\toprule
Defect type & Benchmark & Holdout \\
\midrule
Edge-Ring & 84  & 1{,}302 \\
Edge-Loc  & 53  & 739     \\
Center    & 50  & 603     \\
Loc       & 34  & 492     \\
Scratch   & 15  & 169     \\
Random    &  5  & 108     \\
Donut     &  7  &  71     \\
Near-full &  2  &  16     \\
\bottomrule
\end{tabular}
\end{table}

\begin{table}[H]
\centering
\caption{Top expanded-holdout results among the currently reevaluated leading checkpoints
(70\,k normal / 3.5\,k defect test set).}
\label{tab:holdout_leaderboard}
\begin{tabularx}{\linewidth}{>{\raggedright\arraybackslash}Xcccc}
\toprule
Model / variant & F1 & AUROC & AUPRC & Sweep F1 \\
\midrule
PatchCore + EfficientNet-B1, direct $240{\times}240$ & 0.596 & 0.953 & 0.656 & 0.671 \\
PatchCore + ViT-B/16, direct $224{\times}224$       & 0.548 & 0.941 & 0.615 & 0.606 \\
Teacher-Student + ResNet50                          & 0.516 & 0.926 & 0.622 & 0.585 \\
Autoencoder + BatchNorm (\texttt{max\_abs})        & 0.507 & 0.858 & 0.638 & 0.663 \\
Residual Autoencoder (\texttt{max\_abs})           & 0.503 & 0.871 & 0.637 & 0.643 \\
AE-BN-Dropout-0.00 (\texttt{max\_abs})             & 0.501 & 0.857 & 0.654 & 0.673 \\
\bottomrule
\end{tabularx}
\end{table}

The expanded holdout confirms that the strongest direct PatchCore variants remain
the most stable on a much larger test pool. PatchCore with EfficientNet-B1 now
leads the available holdout reevaluations on deployed F1, AUROC, and AUPRC,
while the ViT-B/16 run remains a strong second and preserves a clear margin over
the teacher-student baseline. The best autoencoder variants still remain
competitive, especially on AUPRC and optimistic threshold-sweep F1, but they no
longer lead the deployed operating point once the test pool expands to
70{,}000 normal and 3{,}500 defect wafers.

This expanded-holdout bundle still does not include the newer WideResNet50-2
$224{\times}224$ PatchCore winner from the main benchmark leaderboard, nor the
fresh FastFlow rerun. It should therefore be read as a robustness check on the
currently reevaluated checkpoints rather than as a complete replacement for the
main benchmark ranking.

\begin{figure}[H]
\centering
\begin{subfigure}[t]{0.48\linewidth}
  \includegraphics[width=\linewidth]{experiments/anomaly_detection/patchcore/efficientnet_b1/x240/main/artifacts/patchcore_efficientnet_b1_one_layer/results/holdout70k_3p5k/plots/holdout_distribution_sweep_confusion.png}
  \caption{PatchCore + EfficientNet-B1 ($\mathrm{F1}=0.596$)}
\end{subfigure}
\hfill
\begin{subfigure}[t]{0.48\linewidth}
  \includegraphics[width=\linewidth]{experiments/anomaly_detection/patchcore/vit_b16/x224/main/artifacts/patchcore_vit_b16_5pct/holdout70k_3p5k/plots/test_evaluation.png}
  \caption{PatchCore + ViT-B/16 ($\mathrm{F1}=0.548$)}
\end{subfigure}
\caption{Expanded-holdout score distributions for the two strongest currently
reevaluated PatchCore checkpoints (70k normal / 3.5k defect test set).
EfficientNet-B1 leads the holdout reevaluations on deployed F1, while ViT-B/16
remains the main-benchmark leader. Both checkpoints use the same
validation-derived threshold from the main benchmark without recalibration.}
\label{fig:holdout_histograms}
\end{figure}


\section{UMAP-Based Embedding Analysis}

UMAP projects high-dimensional wafer embeddings into 2D to reveal whether
normal and anomalous wafers occupy distinct regions of the feature space.
The key question is not whether the UMAP looks clean (wafer defects are
small and sparse, so perfect geometric separation is not expected) but
rather \emph{what pattern of mixing explains the gap between deployed F1 and
best-sweep F1} across all models.

Three figures examine this from different angles. Figure~\ref{fig:umap_comparison}
compares the global baseline against WRN50-2 PatchCore x224 in a binary
split-coloured view. Figure~\ref{fig:umap_vit_patchcore} contrasts the
split-coloured and score-coloured views for the best model (ViT-B/16 PatchCore),
showing that even when the geometry looks mixed, the model correctly concentrates
anomaly scores where defects actually are. Figure~\ref{fig:umap_defect_types}
provides the most informative view: the full backbone progression from baseline
to ViT, coloured by defect type rather than binary normal/anomaly.

\begin{figure}[H]
\centering
\begin{subfigure}[t]{0.48\linewidth}
  \includegraphics[width=\linewidth]{experiments/anomaly_detection/backbone_embedding/wide_resnet50_2/x64/baseline/artifacts/umaps/wideresnet50A_embedding_baseline/evaluation/plots/embedding_umap.png}
  \caption{WRN50-2 center-distance baseline: anomalies (red) are distributed
  uniformly throughout the same islands as normals. No geometric threshold can
  reliably separate them.}
  \label{fig:umap_wrn_baseline}
\end{subfigure}
\hfill
\begin{subfigure}[t]{0.48\linewidth}
  \includegraphics[width=\linewidth]{experiments/anomaly_detection/patchcore/wideresnet50/x224/multilayer_umap_followup/artifacts/patchcore-wideresnet50-multilayer-umap/topk_mb50k_r005_x224/plots/umap_by_split.png}
  \caption{WRN50-2 PatchCore $224\!\times\!224$: anomalies begin to concentrate
  on manifold boundaries, thinner branches, and side islands rather than the
  dominant normal core, healthier geometry than the baseline, but still mixed.}
  \label{fig:umap_wrn_patchcore}
\end{subfigure}
\caption{Split-coloured UMAP comparison: global center-distance baseline (left)
vs. the WideResNet50-2 PatchCore $224\!\times\!224$ checkpoint (right).
Anomaly wafers are shown in red, normals in blue/grey. The progression from
completely uniform mixing to boundary-concentrated anomalies reflects the gain
from switching to local patch-level scoring.}
\label{fig:umap_comparison}
\end{figure}

\begin{figure}[H]
\centering
\begin{subfigure}[t]{0.48\linewidth}
  \includegraphics[width=\linewidth]{experiments/anomaly_detection/patchcore/vit_b16/x224/main/artifacts/patchcore_vit_b16_5pct/main_5pct/plots/umap_test_embeddings.png}
  \caption{Split-coloured: anomalies remain mixed within the normal manifold,
  explaining why a simple geometric threshold cannot close the gap between
  deployed and best-sweep F1.}
\end{subfigure}
\hfill
\begin{subfigure}[t]{0.48\linewidth}
  \includegraphics[width=\linewidth]{experiments/anomaly_detection/patchcore/vit_b16/x224/main/artifacts/patchcore_vit_b16_5pct/main_5pct/plots/umap_by_score.png}
  \caption{Score-coloured: the ViT PatchCore score concentrates high anomaly
  evidence in specific compact regions and tails, confirming the model has
  learned strong anomaly signal even though the 2D geometry looks mixed.}
\end{subfigure}
\caption{ViT-B/16 PatchCore $224\!\times\!224$ UMAP (cosine metric, fit on
5,000 train-normal reference points). The contrast between the two panels is the
key diagnostic: the split view shows \emph{why} a UMAP-space KNN threshold fails
(F1 = 0.083), while the score view confirms \emph{that} the model correctly
assigns high scores to anomaly-concentrated regions, giving F1 = 0.595 on the
original wafer score.}
\label{fig:umap_vit_patchcore}
\end{figure}

\begin{figure}[H]
\centering
\includegraphics[width=\linewidth]{experiments/anomaly_detection/report_figures/artifacts/plots/umap_defect_type_comparison.png}
\caption{UMAP coloured by defect type across the backbone progression (left to right:
WRN50-2 center-distance baseline, ResNet18 PatchCore x64, EfficientNet-B1 PatchCore
x240, ViT-B/16 PatchCore x224). Top row: each defect type in a distinct colour,
normals in grey. Bottom row: continuous anomaly score heatmap.
In the baseline (leftmost), anomaly points are uniformly distributed throughout
the normal cloud with no geometric separation. As backbone quality improves,
anomalies begin to concentrate on manifold boundaries and side islands (middle),
and the score-coloured view shows increasingly concentrated anomaly evidence
even where the split view still looks mixed. Scratch points (brown, few samples)
remain the hardest to separate geometrically across all models.
Generated by \texttt{experiments/anomaly\_detection/report\_figures/notebook.ipynb}, Plot 6.}
\label{fig:umap_defect_types}
\end{figure}

\textbf{What the UMAPs reveal.} The overlapping geometry is not a model failure
--- it is the defining characteristic of the wafer defect detection problem.
Defects are spatially small and sparse; a defective wafer may differ from a normal
one in only a few patch locations, leaving its global embedding largely inside the
normal manifold. This is why global center-distance scoring (left panel of
Figure~\ref{fig:umap_comparison}) cannot work: there is simply no useful global
separation to threshold. Local patch-level methods (PatchCore, ViT) succeed
\emph{despite} this mixing because they score each patch independently rather
than relying on a global embedding distance. The score-coloured ViT UMAP
(Figure~\ref{fig:umap_vit_patchcore}, right) shows precisely this: even though
anomaly wafers are geometrically mixed with normals, the model knows \emph{where}
to concentrate anomaly evidence, producing a strong wafer-level score. The
remaining gap between deployed F1 (0.595) and best-sweep F1 (0.651) is a
threshold-calibration problem, not a feature-quality problem. A more
principled threshold rule operating on the score distribution directly would
recover most of this headroom.

\section{Thresholding Limitations and Discussion}

The 95th-percentile validation-normal threshold is fair, reproducible, and
requires no defect labels.  However, three consistent observations from the
diagnostics show its limitations:

\begin{itemize}[leftmargin=1.5em]
  \item \textbf{Score overlap.} Even the best models show substantial overlap
        between normal and anomaly score distributions.  The histogram for
        the WRN PatchCore x224 selected variant (Figure~\ref{fig:wrn_x224_diag})
        shows that the normal and anomaly distributions overlap near the
        operating threshold.
  \item \textbf{Geometric overlap.} UMAP figures show that anomalies are not
        cleanly peeled off from the normal manifold. They occupy boundaries
        and tails, but do not form an isolated cluster.
  \item \textbf{Sweep headroom.} Best-sweep F1 is systematically $0.06$--$0.08$
        above deployed-threshold F1 across the strongest models (e.g.\ WRN
        PatchCore x224: deployed 0.549, sweep 0.634).  This gap represents
        the threshold-selection headroom that a more principled rule could capture.
\end{itemize}

Possible directions to improve thresholding beyond the 95th-percentile rule:
\begin{itemize}[leftmargin=1.5em]
  \item Adaptive thresholding calibrated on validation score structure.
  \item Local-density or manifold-aware decision boundaries informed by embedding geometry.
        The current UMAP-space KNN follow-ups are not yet sufficient:
        for example, the ViT-B/16 branch keeps a very strong original wafer
        score while its 2-D UMAP-KNN score collapses to near-random ranking.
  \item Hybrid supervision: if any defect examples can be used for threshold tuning
        without contaminating training, they could inform a more calibrated operating point.
\end{itemize}

Not every main-report winner has been reevaluated on the expanded holdout yet.
The holdout bundle is therefore a stability check on the reevaluated subset,
not a universal replacement benchmark.

\section{Conclusion}

This project demonstrates a clear and reproducible progression in wafer-map
anomaly detection, implemented entirely in PyTorch:

\begin{enumerate}[leftmargin=1.5em]
  \item Plain reconstruction models establish a viable pipeline and show that
        local-defect sensitivity is the main limitation.
  \item Better in-domain autoencoder variants and score selection improve
        performance without new training. The BatchNorm autoencoder scored with
        maximum absolute pixel error is the strongest reconstruction baseline
        ($\mathrm{F1}=0.502$).
  \item Pretrained backbones with global center-distance scoring are consistently
        weak regardless of backbone size.
  \item Local anomaly scoring methods (teacher-student distillation and
        PatchCore) produce the strongest results. The score composition
        choice interacts strongly with the checkpoint. Post-training score
        sweeps are as important as architecture choices.
  \item Native high-resolution preprocessing provides a consistent and material
        lift: EfficientNet-B0 gained $+0.077$ F1 from $64\!\times\!64$ to
        $224\!\times\!224$. EfficientNet-B1 at $240\!\times\!240$ reaches
        $\mathrm{F1}=0.591$, the strongest CNN result in the project.
  \item The best completed result on the main benchmark, ViT-B/16 PatchCore at
        direct $224\!\times\!224$, achieves $\mathrm{F1}=0.595$,
        $\mathrm{AUROC}=0.956$, and $\mathrm{AUPRC}=0.671$. It leads
        EfficientNet-B1 by only 0.004 F1, with the ViT's advantage concentrated
        in Scratch, Loc, and Edge-Loc recall.
\end{enumerate}

The dual conclusion is: representation quality and local scoring matter
greatly, and the combination of a strong pretrained backbone with patch-level
nearest-neighbour scoring substantially outperforms reconstruction-based
or globally aggregated alternatives.  At the same time, the remaining
bottleneck is now score overlap and threshold selection rather than feature
extraction, which points toward more principled operating-point selection as
the highest-leverage next step.

\appendix
\newpage
\section{Appendix: Group Contributions}
\label{app:contributions}

\begin{tabularx}{\linewidth}{>{\raggedright\arraybackslash}p{3.5cm}>{\raggedright\arraybackslash}X}
\toprule
Member & Primary contributions \\
\midrule
Henry Lee Jun (1004219)
  & Autoencoder family (baseline, BatchNorm, dropout sweep, residual, 128px variant),
    score ablation framework, Deep SVDD, VAE beta sweep, teacher-student distillation
    (ResNet18 and ResNet50), WideResNet50-2 teacher-student variants, score ensemble
    study, report writing and LaTeX preparation. \\
Chia Tang (1007200)
  &  \\
Genson Low (1005931)
  & \\
\bottomrule
\end{tabularx}

\newpage
\section{Appendix: How to Reproduce Each Figure}
\label{app:figures}

The rubric requires that each figure in the report can be recreated independently.
The table below maps every figure to the notebook or script that generates it.
All notebooks load pre-trained checkpoints by default (\texttt{RETRAIN = False})
and regenerate all plots from saved artifacts without requiring GPU access.
Setting \texttt{RETRAIN = True} reruns the full training and evaluation pipeline.

\begin{tabularx}{\linewidth}{>{\raggedright\arraybackslash}p{0.38\linewidth}>{\raggedright\arraybackslash}X}
\toprule
Figure & Source notebook \\
\midrule
Defect gallery (Fig.~\ref{fig:defect_gallery})
  & \texttt{artifacts/report\_plots/defect\_type\_gallery.png}, generated by the dataset notebook at \texttt{data/dataset/x64/benchmark\_50k\_5pct/notebook.ipynb} \\
Overall comparison (Fig.~\ref{fig:overall_comparison})
  & \texttt{artifacts/report\_plots/overall\_experiment\_comparison.png}, generated by running the report-plot aggregation cell in any top-level analysis notebook \\
AE training and score ablation (Fig.~\ref{fig:ae_training_ablation})
  & \texttt{experiments/anomaly\_detection/autoencoder/x64/baseline/notebook.ipynb} \\
AE x224 score ablation (Fig.~\ref{fig:ae224_score_ablation})
  & \texttt{experiments/anomaly\_detection/autoencoder/x224/main/artifacts/autoencoder\_x224/plots/score\_ablation\_summary.png}, generated by the x224 notebook score ablation cell \\
AE reconstruction examples (Fig.~\ref{fig:ae_reconstruction})
  & \texttt{experiments/anomaly\_detection/autoencoder/x64/baseline/notebook.ipynb} \\
AE family comparison (Fig.~\ref{fig:ae_family})
  & \texttt{artifacts/report\_plots/autoencoder\_family\_comparison.png}, regenerated by the autoencoder baseline notebook after all variant artifacts are present \\
Compact baselines (Fig.~\ref{fig:compact_baselines})
  & \texttt{artifacts/report\_plots/compact\_baseline\_comparison.png}, requires VAE, SVDD, and backbone-embedding artifacts \\
Teacher-student family (Fig.~\ref{fig:ts_family})
  & \texttt{artifacts/report\_plots/ts\_family\_comparison.png}, requires ResNet18 and ResNet50 TS artifacts \\
FastFlow results (Fig.~\ref{fig:fastflow_results})
  & \texttt{experiments/anomaly\_detection/fastflow/x64/main/notebook.ipynb} \\
PatchCore family comparison (Fig.~\ref{fig:patchcore_family})
  & \texttt{artifacts/report\_plots/patchcore\_family\_comparison.png}, requires all PatchCore variant artifacts \\
EfficientNet-B1 results (Fig.~\ref{fig:effnetb1_x240})
  & \texttt{experiments/anomaly\_detection/patchcore/efficientnet\_b1/x240/main/notebook.ipynb} \\
WRN50-2 x224 diagnostics (Fig.~\ref{fig:wrn_x224_diag})
  & \texttt{experiments/anomaly\_detection/patchcore/wideresnet50/x224/multilayer/notebook.ipynb} \\
ViT-B/16 main result (Fig.~\ref{fig:vit_main_result})
  & \texttt{experiments/anomaly\_detection/patchcore/vit\_b16/x224/main/notebook.ipynb} \\
Holdout histograms (Fig.~\ref{fig:holdout_histograms})
  & EfficientNet-B1 and ViT holdout sections in their respective notebooks above. TS-ResNet50 from \texttt{experiments/anomaly\_detection/teacher\_student/resnet50/x64/main/notebook.ipynb} \\
WRN baseline UMAP (Fig.~\ref{fig:umap_wrn_baseline})
  & \texttt{experiments/anomaly\_detection/backbone\_embedding/wide\_resnet50\_2/x64/baseline/notebook.ipynb} \\
ResNet18 PatchCore UMAP (Fig.~\ref{fig:umap_resnet18_patchcore})
  & \texttt{experiments/anomaly\_detection/patchcore/resnet18/x64/holdout70k\_3p5k\_umap\_followup/notebook.ipynb} \\
WRN50-2 PatchCore x224 UMAP (Fig.~\ref{fig:umap_wrn_patchcore})
  & \texttt{experiments/anomaly\_detection/patchcore/wideresnet50/x224/multilayer\_umap\_followup/notebook.ipynb} \\
ViT PatchCore UMAP (Fig.~\ref{fig:umap_vit_patchcore})
  & \texttt{experiments/anomaly\_detection/patchcore/vit\_b16/x224/main/notebook.ipynb} (UMAP section) \\
\midrule
\multicolumn{2}{l}{\textit{Cross-model comparison figures -- all generated by
\texttt{experiments/anomaly\_detection/report\_figures/notebook.ipynb}}} \\
Per-defect recall chart (Fig.~\ref{fig:per_defect_chart})
  & \texttt{report\_figures/notebook.ipynb}, Plot 1 \\
Resolution effect comparison (Fig.~\ref{fig:resolution_comparison})
  & \texttt{report\_figures/notebook.ipynb}, Plot 2 \\
PatchCore backbone progression (Fig.~\ref{fig:patchcore_progression})
  & \texttt{report\_figures/notebook.ipynb}, Plot 3 \\
Score distribution overlay (Fig.~\ref{fig:score_distributions})
  & \texttt{report\_figures/notebook.ipynb}, Plot 4 \\
TS family full comparison (Fig.~\ref{fig:ts_family_full})
  & \texttt{report\_figures/notebook.ipynb}, Plot 5 \\
Improved UMAP by defect type (Fig.~\ref{fig:umap_defect_types})
  & \texttt{report\_figures/notebook.ipynb}, Plot 6 \\
EfficientNet-B0 resolution UMAP (Fig.~\ref{fig:umap_resolution_effb0})
  & Feature extraction script run via \texttt{report\_figures/notebook.ipynb}; CSVs saved to \texttt{experiments/anomaly\_detection/patchcore/efficientnet\_b0/x64/umap\_points.csv} and \texttt{.../x224/umap\_points.csv} \\
\bottomrule
\end{tabularx}

\newpage
\section{Appendix: Experiment Inventory}
\label{app:notebooks}

This appendix documents the broader search space behind the main narrative.
Table~\ref{tab:all_experiments} lists every training run completed in this
project. Subsequent tables break down sweep results within each family.

\subsection{Complete Experiment List}

\begin{longtable}{>{\raggedright\arraybackslash}p{4.8cm}
                  >{\raggedright\arraybackslash}p{1.4cm}
                  c c c}
\caption{All experiments run in this project (main benchmark, best selected score per run).
Runs marked with $\dagger$ are intermediate steps superseded by a later variant in the same family.}
\label{tab:all_experiments} \\
\toprule
Experiment & Resolution & F1 & AUROC & AUPRC \\
\midrule
\endfirsthead
\multicolumn{5}{c}{\textit{(continued from previous page)}} \\
\toprule
Experiment & Resolution & F1 & AUROC & AUPRC \\
\midrule
\endhead
\midrule
\multicolumn{5}{r}{\textit{continued on next page}} \\
\endfoot
\bottomrule
\endlastfoot

\multicolumn{5}{l}{\textit{Reconstruction family}} \\
AE baseline (\texttt{topk\_abs\_mean}) & 64$\times$64 & 0.468 & 0.839 & 0.522 \\
AE baseline, longer run (43 ep)$^\dagger$ & 64$\times$64 & 0.460 & 0.835 & 0.525 \\
AE + BatchNorm (\texttt{max\_abs}) & 64$\times$64 & 0.502 & 0.834 & 0.568 \\
AE + BatchNorm + Dropout 0.00 & 64$\times$64 & 0.487 & 0.851 & 0.617 \\
AE + BatchNorm + Dropout 0.05 & 64$\times$64 & 0.483 & 0.835 & 0.552 \\
AE + BatchNorm + Dropout 0.10 & 64$\times$64 & 0.484 & 0.845 & 0.570 \\
AE + BatchNorm + Dropout 0.20 & 64$\times$64 & 0.470 & 0.841 & 0.575 \\
Residual AE (\texttt{max\_abs}) & 64$\times$64 & 0.474 & 0.843 & 0.589 \\
AE baseline (\texttt{topk\_abs\_mean}) & 128$\times$128 & 0.431 & 0.815 & 0.455 \\
VAE ($\beta=0.001$) & 64$\times$64 & 0.343 & 0.770 & 0.336 \\
VAE ($\beta=0.005$) & 64$\times$64 & 0.340 & 0.772 & 0.372 \\
VAE ($\beta=0.01$) & 64$\times$64 & 0.335 & 0.766 & 0.369 \\
VAE ($\beta=0.05$) & 64$\times$64 & 0.302 & 0.750 & 0.329 \\
Deep SVDD & 64$\times$64 & 0.360 & 0.788 & 0.213 \\
\midrule

\multicolumn{5}{l}{\textit{Pretrained backbone center-distance baselines}} \\
ResNet18 center-distance & 64$\times$64 & 0.236 & 0.685 & 0.195 \\
WideResNet50-2 center-distance & 64$\times$64 & 0.243 & 0.677 & 0.142 \\
\midrule

\multicolumn{5}{l}{\textit{Teacher-student distillation}} \\
TS + ResNet18 (student-only topk20) & 64$\times$64 & 0.495 & 0.894 & 0.519 \\
TS + ResNet50 (mixed 2:1 topk20) & 64$\times$64 & 0.525 & 0.909 & 0.599 \\
TS + WideResNet50-2 Layer2 (topk25) & 64$\times$64 & 0.508 & 0.920 & 0.540 \\
TS + WideResNet50-2 Multilayer (topk15) & 64$\times$64 & 0.524 & 0.923 & 0.546 \\
\midrule

\multicolumn{5}{l}{\textit{FastFlow}} \\
FastFlow WRN50-2 layer2+layer3, 4 steps & 64$\times$64 & 0.482 & 0.871 & 0.489 \\
FastFlow WRN50-2 layer2+layer3, 6 steps$^\dagger$ & 64$\times$64 & 0.470 & 0.870 & 0.479 \\
FastFlow WRN50-2 layer2 only, 6 steps$^\dagger$ & 64$\times$64 & 0.442 & 0.884 & 0.460 \\
\midrule

\multicolumn{5}{l}{\textit{PatchCore}} \\
PatchCore + AE-BN encoder (\texttt{mean}) & 64$\times$64 & 0.336 & 0.851 & 0.226 \\
PatchCore + ResNet18 (\texttt{mean}) & 64$\times$64 & 0.401 & 0.842 & 0.411 \\
PatchCore + ResNet50 (\texttt{mean}) & 64$\times$64 & 0.420 & 0.821 & 0.363 \\
PatchCore + WideResNet50-2 (\texttt{topk r=0.10}) & 64$\times$64 & 0.532 & 0.917 & 0.562 \\
PatchCore + WideResNet50-2 (\texttt{topk r=0.05}) & 64$\times$64 & 0.525 & 0.921 & 0.564 \\
PatchCore + WideResNet50-2 (\texttt{topk r=0.15}) & 64$\times$64 & 0.526 & 0.912 & 0.549 \\
PatchCore + EfficientNet-B0 (\texttt{topk r=0.02})$^\dagger$ & 64$\times$64 & 0.467 & 0.905 & 0.489 \\
PatchCore + WideResNet50-2 (\texttt{topk r=0.05}) & 224$\times$224 & 0.549 & 0.931 & 0.659 \\
PatchCore + EfficientNet-B0 (\texttt{topk r=0.02}) & 224$\times$224 & 0.544 & 0.925 & 0.483 \\
PatchCore + EfficientNet-B1 (\texttt{topk r=0.03}) & 240$\times$240 & 0.591 & 0.935 & 0.609 \\
PatchCore + ViT-B/16 block 6 (\texttt{topk r=0.10}) & 224$\times$224 & \textbf{0.595} & \textbf{0.956} & \textbf{0.671} \\
\midrule

\multicolumn{5}{l}{\textit{Score ensemble}} \\
WRN PatchCore + TS-Res50 (max fusion) & 64$\times$64 & 0.529 & 0.916 & 0.611 \\

\end{longtable}

\subsection{Autoencoder-Family Leaderboard (Main Benchmark)}

\begin{table}[H]
\centering
\caption{Selected autoencoder-family results (main benchmark, best score per checkpoint).}
\label{tab:appendix_ae}
\begin{tabular}{lccc}
\toprule
Variant & F1 & AUROC & AUPRC \\
\midrule
AE-224-topk               & 0.510 & 0.901 & 0.596 \\
AE-64-BN-max              & 0.502 & 0.834 & 0.568 \\
AE-64-BN-DO0.00-max       & 0.487 & 0.851 & 0.617 \\
AE-64-BN-DO0.10-max       & 0.484 & 0.845 & 0.570 \\
AE-64-BN-DO0.05-max       & 0.483 & 0.835 & 0.552 \\
AE-64-Res-max             & 0.474 & 0.843 & 0.589 \\
AE-64-topk (43 ep)        & 0.460 & 0.835 & 0.525 \\
AE-64-BN-topk             & 0.431 & 0.790 & 0.603 \\
AE-128-topk               & 0.431 & 0.815 & 0.455 \\
\bottomrule
\end{tabular}
\end{table}

\subsection{PatchCore Leaderboard (Main Benchmark)}

\begin{table}[H]
\centering
\caption{Selected PatchCore results (main benchmark, best variant per backbone).}
\label{tab:appendix_patchcore}
\begin{tabular}{lccc}
\toprule
Variant & F1 & AUROC & AUPRC \\
\midrule
ViT-B16-x224-topk10-mb400k        & 0.595 & 0.956 & 0.671 \\
EffNetB1-x240-topk3-mb240k        & 0.591 & 0.935 & 0.609 \\
WideRes50-x224-topk-mb600k-r005   & 0.549 & 0.931 & 0.659 \\
EffNetB0-x224-topk-mb240k-r002    & 0.544 & 0.925 & 0.483 \\
WideRes50-topk-mb50k-r010         & 0.532 & 0.917 & 0.562 \\
WideRes50-topk-mb50k-r015         & 0.526 & 0.912 & 0.549 \\
WideRes50-topk-mb50k-r005         & 0.525 & 0.921 & 0.564 \\
EffNetB0-x64-topk-mb240k-r002     & 0.467 & 0.905 & 0.489 \\
Res50-mean-mb50k                  & 0.420 & 0.821 & 0.363 \\
Res18-mean-mb50k                  & 0.401 & 0.842 & 0.411 \\
AE-BN-mean-mb50k                  & 0.336 & 0.851 & 0.226 \\
\bottomrule
\end{tabular}
\end{table}

\subsection{AE-Backed PatchCore Sweep}

\begin{table}[H]
\centering
\caption{PatchCore variants built on the frozen AE-BN encoder.}
\label{tab:appendix_aebn_patchcore}
\begin{tabularx}{\linewidth}{lcc>{\raggedright\arraybackslash}X}
\toprule
Variant & F1 & AUROC & Notes \\
\midrule
\texttt{mean\_mb50k}       & 0.336 & 0.851 & Best AE-BN PatchCore result \\
\texttt{topk\_mb50k\_r010} & 0.185 & 0.809 & Top-k reduction, 50\,k bank \\
\texttt{topk\_mb50k\_r005} & 0.123 & 0.777 & Narrower top-k region \\
\texttt{max\_mb50k}        & 0.056 & 0.679 & Max reduction too brittle \\
\bottomrule
\end{tabularx}
\end{table}

\subsection{ResNet18 Teacher-Student Focused Ablation}

\begin{table}[H]
\centering
\caption{ResNet18 teacher-student focused layer and top-k ablation.}
\label{tab:appendix_ts_resnet18}
\begin{tabular}{lcccc}
\toprule
Variant & Layer & Top-k ratio & F1 & AUROC \\
\midrule
\texttt{ts\_resnet18\_layer2\_topk25} & layer2 & 0.25 & 0.494 & 0.890 \\
\texttt{ts\_resnet18\_layer2\_topk20} & layer2 & 0.20 & 0.494 & 0.889 \\
\texttt{ts\_resnet18\_layer2\_topk15} & layer2 & 0.15 & 0.483 & 0.889 \\
\texttt{ts\_resnet18\_layer1\_topk15} & layer1 & 0.15 & 0.411 & 0.824 \\
\texttt{ts\_resnet18\_layer1\_topk20} & layer1 & 0.20 & 0.387 & 0.817 \\
\bottomrule
\end{tabular}
\end{table}

\subsection{ResNet50 Teacher-Student Score Sweep}

\begin{table}[H]
\centering
\caption{ResNet50 teacher-student local score-sweep summary (top rows).}
\label{tab:appendix_ts_resnet50}
\begin{tabular}{lcccc}
\toprule
Score variant & F1 & AUROC & AUPRC & Sweep F1 \\
\midrule
\texttt{s2\_a1\_topk\_mean\_r0.20} (selected) & 0.525 & 0.909 & 0.599 & 0.606 \\
\texttt{s1\_a0.5\_topk\_mean\_r0.20}           & 0.525 & 0.909 & 0.599 & 0.606 \\
\texttt{s1\_a2\_topk\_mean\_r0.20}             & 0.523 & 0.904 & 0.570 & 0.565 \\
\texttt{s2\_a1\_topk\_mean\_r0.10}             & 0.506 & 0.913 & 0.584 & 0.563 \\
\bottomrule
\end{tabular}
\end{table}

\subsection{WideResNet50-2 Teacher-Student Score Sweep}

\begin{table}[H]
\centering
\caption{WideResNet50-2 multilayer teacher-student score-sweep (top rows).}
\label{tab:appendix_ts_wrn}
\begin{tabular}{lcccc}
\toprule
Score variant & F1 & AUROC & AUPRC & Sweep F1 \\
\midrule
\texttt{s2\_a1\_topk\_mean\_r0.15} (selected) & 0.524 & 0.923 & 0.546 & 0.561 \\
\texttt{s1\_a1\_topk\_mean\_r0.25}             & 0.518 & 0.926 & 0.538 & 0.558 \\
\texttt{s2\_a1\_topk\_mean\_r0.20}             & 0.512 & 0.922 & 0.543 & 0.561 \\
\texttt{s2\_a1\_topk\_mean\_r0.10}             & 0.511 & 0.923 & 0.544 & 0.556 \\
\bottomrule
\end{tabular}
\end{table}

\subsection{FastFlow Ablation Summary}

\begin{table}[H]
\centering
\caption{FastFlow training-variant ablation (WideResNet50-2 backbone).}
\label{tab:appendix_fastflow}
\begin{tabular}{lccccc}
\toprule
Training variant & Layers & Steps & Best score & F1 & AUROC \\
\midrule
\texttt{wrn50\_l23\_s4} (selected) & layer2+layer3 & 4 & mean       & 0.482 & 0.871 \\
\texttt{wrn50\_l23\_s6}            & layer2+layer3 & 6 & mean       & 0.470 & 0.870 \\
\texttt{wrn50\_l2\_s6}             & layer2 only   & 6 & topk-r0.15 & 0.442 & 0.884 \\
\bottomrule
\end{tabular}
\end{table}

\subsection{Score Ensemble Study}

\begin{table}[H]
\centering
\caption{Post-hoc score ensemble: WRN PatchCore + TS-Res50.}
\label{tab:appendix_ensemble}
\begin{tabular}{lcccc}
\toprule
Fusion variant & F1 & AUROC & AUPRC & Sweep F1 \\
\midrule
PatchCore only         & 0.532 & 0.917 & 0.562 & 0.586 \\
max pair (selected)    & 0.529 & 0.916 & \textbf{0.611} & 0.620 \\
mean 70/30 PatchCore   & 0.528 & 0.918 & 0.596 & 0.616 \\
mean 50/50             & 0.526 & 0.917 & 0.606 & 0.615 \\
TS only                & 0.525 & 0.909 & 0.599 & 0.606 \\
\bottomrule
\end{tabular}
\end{table}

The ensemble improved AUPRC notably (0.611 vs 0.562) but did not beat the
standalone PatchCore on deployed F1.

\subsection{Experiment-to-File Reference}

\begin{table}[H]
\centering
\caption{Key experiment-to-notebook mapping for reproducibility.
All paths are relative to the project root.}
\label{tab:appendix_notebooks}
\begin{tabularx}{\linewidth}{>{\raggedright\arraybackslash}X>{\raggedright\arraybackslash}X}
\toprule
Experiment & Notebook path \\
\midrule
Baseline AE (224$\times$224)
  & \texttt{experiments/anomaly\_detection/autoencoder/x224/main/notebook.ipynb} \\
Baseline AE
  & \texttt{experiments/anomaly\_detection/autoencoder/x64/baseline/notebook.ipynb} \\
BatchNorm AE
  & \texttt{experiments/anomaly\_detection/autoencoder/x64/batchnorm/notebook.ipynb} \\
BatchNorm + Dropout sweep
  & \texttt{experiments/anomaly\_detection/autoencoder/x64/batchnorm\_dropout/notebook.ipynb} \\
Residual AE
  & \texttt{experiments/anomaly\_detection/autoencoder/x64/residual/notebook.ipynb} \\
AE 128$\times$128
  & \texttt{experiments/anomaly\_detection/autoencoder/x128/baseline/notebook.ipynb} \\
VAE baseline
  & \texttt{experiments/anomaly\_detection/vae/x64/baseline/notebook.ipynb} \\
VAE beta sweep
  & \texttt{experiments/anomaly\_detection/vae/x64/beta\_sweep/notebook.ipynb} \\
Deep SVDD
  & \texttt{experiments/anomaly\_detection/svdd/x64/baseline/notebook.ipynb} \\
ResNet18 backbone baseline
  & \texttt{experiments/anomaly\_detection/backbone\_embedding/resnet18/x64/baseline/notebook.ipynb} \\
WRN50-2 embedding baseline
  & \texttt{experiments/anomaly\_detection/backbone\_embedding/wide\_resnet50\_2/x64/baseline/notebook.ipynb} \\
PatchCore + AE-BN backbone
  & \texttt{experiments/anomaly\_detection/patchcore/ae\_bn/x64/main/notebook.ipynb} \\
PatchCore + ResNet18
  & \texttt{experiments/anomaly\_detection/patchcore/resnet18/x64/main/notebook.ipynb} \\
PatchCore + ResNet50
  & \texttt{experiments/anomaly\_detection/patchcore/resnet50/x64/main/notebook.ipynb} \\
PatchCore + WRN50-2 x64
  & \texttt{experiments/anomaly\_detection/patchcore/wideresnet50/x64/main/notebook.ipynb} \\
PatchCore + WRN50-2 x224
  & \texttt{experiments/anomaly\_detection/patchcore/wideresnet50/x224/multilayer/notebook.ipynb} \\
PatchCore + WRN50-2 x224 UMAP
  & \texttt{experiments/anomaly\_detection/patchcore/wideresnet50/x224/multilayer\_umap\_followup/notebook.ipynb} \\
PatchCore + EfficientNet-B0 x224
  & \texttt{experiments/anomaly\_detection/patchcore/efficientnet\_b0/x224/main/notebook.ipynb} \\
PatchCore + EfficientNet-B1 x240
  & \texttt{experiments/anomaly\_detection/patchcore/efficientnet\_b1/x240/main/notebook.ipynb} \\
PatchCore + ViT-B/16 x224
  & \texttt{experiments/anomaly\_detection/patchcore/vit\_b16/x224/main/notebook.ipynb} \\
TS + ResNet18
  & \texttt{experiments/anomaly\_detection/teacher\_student/resnet18/x64/main/notebook.ipynb} \\
TS + ResNet50
  & \texttt{experiments/anomaly\_detection/teacher\_student/resnet50/x64/main/notebook.ipynb} \\
TS + WRN50-2 layer2
  & \texttt{experiments/anomaly\_detection/teacher\_student/wideresnet50\_2/x64/layer2\_self\_contained/notebook.ipynb} \\
TS + WRN50-2 multilayer
  & \texttt{experiments/anomaly\_detection/teacher\_student/wideresnet50\_2/x64/multilayer\_self\_contained/notebook.ipynb} \\
FastFlow sweep
  & \texttt{experiments/anomaly\_detection/fastflow/x64/main/notebook.ipynb} \\
Score ensemble
  & \texttt{experiments/anomaly\_detection/ensemble/x64/score\_ensemble/notebook.ipynb} \\
Report comparison figures
  & \texttt{experiments/anomaly\_detection/report\_figures/notebook.ipynb} \\
\bottomrule
\end{tabularx}
\end{table}

\newpage
\begin{thebibliography}{9}

\bibitem{khatun2026cbam}
M.~R. Khatun, F.~A. Farid, S.~Dhar, M.~S. Islam, J.~Uddin, and H.~A. Karim,
\newblock ``CBAM-enhanced lightweight CNN for wafer map defect classification,''
\newblock \textit{Frontiers in Electronics}, vol.~7, 2026.
\newblock \url{https://doi.org/10.3389/felec.2026.1750707}

\bibitem{taha2025survey}
K.~Taha,
\newblock ``Observational and experimental insights into machine learning-based defect classification in wafers,''
\newblock \textit{Journal of Intelligent Manufacturing}, 2025.
\newblock \url{https://doi.org/10.1007/s10845-024-02521-0}

\bibitem{jo2021oneclass}
H.~Y. Jo and S.-W. Lee,
\newblock ``One-class classification for wafer map using adversarial autoencoder with DSVDD prior,''
\newblock \textit{arXiv preprint arXiv:2107.08823}, 2021.
\newblock \url{https://arxiv.org/abs/2107.08823}

\bibitem{roth2022patchcore}
K.~Roth, L.~Pemula, J.~Zepeda, B.~Sch\"{o}lkopf, T.~Brox, and P.~Gehler,
\newblock ``Towards total recall in industrial anomaly detection,''
\newblock in \textit{Proceedings of the IEEE/CVF Conference on Computer Vision and Pattern Recognition (CVPR)},
\newblock pp.\ 14318--14328, 2022.
\newblock \url{https://arxiv.org/abs/2106.08265}

\bibitem{ruff2018deepsvdd}
L.~Ruff, R.~Vandermeulen, N.~Goernitz, L.~Deecke, S.~A. Siddiqui, A.~Binder, E.~M\"{u}ller, and M.~Kloft,
\newblock ``Deep one-class classification,''
\newblock in \textit{Proceedings of the 35th International Conference on Machine Learning (ICML)},
\newblock vol.\ 80, pp.\ 4393--4402, 2018.
\newblock \url{http://proceedings.mlr.press/v80/ruff18a.html}

\end{thebibliography}

\end{document}